\section{No local certificate: the obstruction theorem}\label{sec:obstruction}

A natural strategy for proving \eqref{eq:DOM} is a \emph{certificate}: an identity or inequality
$\Q(f_0s)\ge\int W(x)\,|(Ss)(x)|^2dx$ with an explicit nonnegative weight $W$ and an explicit finite difference
stencil $S$, possibly summed over a family of stencils, chosen so that the right-hand side visibly pays the negative
part of $\nu$. On finite graphs this is exactly how positive semidefinite matrices with zero row sums are
certified (sums of squares of zero-mean stencils). The one positive minorant of the Weil form that is known unconditionally is
nonlocal: for $g$ supported in $[2^{-1/2},2^{1/2}]$ with $\widehat g$ vanishing at $2i$ and $0$, compression to Sonin's space gives
$W_\infty(g*g^*)\ge\operatorname{Tr}(\vartheta(g)S\vartheta(g)^*)$ \cite[Thm.~7.1]{Connes2026}. The theorem of this section says
that no such minorant can be local: every nonnegative finite-stencil minorant of $\Q(f_0\,\cdot)$ is zero. The reason is not
the negative density but the radical, through the linear independence of translates, a special case of the independence question
raised by Bombieri \cite[Introduction, (iii)]{BOMBIERI-WEILQF-2000}.

\begin{lemma}[compact witnesses]\label{lem:budget}
Fix smooth cutoffs $0\le\chi_R\le1$, $\chi_R=1$ on $[-R,R]$, $\chi_R=0$ off $[-R-1,R+1]$, $|\chi_R'|\le c_\chi$.
For $q\in\R$ and $R\ge1$ let $r_q=f_0(\cdot-q)/f_0$ and $s_{q,R}=\chi_R\,r_q\in C_c^\infty$. Then, with $a_q=\frac\pi2e^{-2|q|}$,
\begin{equation}\label{eq:budget}
 |\Q(f_0s_{q,R})|\le C_{\Xf}\,C_q^2\,e^{-2a_qe^{2R}},\qquad
 C_q^2=\frac{\frac{35}{3}A_0^2+\frac43(A_1^2+c_\chi^2A_0^2)}{2a_q}.
\end{equation}
\end{lemma}
\begin{proof}
$f_0s_{q,R}=U_qf_0-e_R$ with $e_R=(1-\chi_R)U_qf_0$. By \Cref{prop:radical}, $\Q(U_qf_0-e_R)=\Q(e_R)$, and
$|\Q(e_R)|\le C_{\Xf}\|e_R\|_{\Xf}^2$. The envelope of \Cref{thm:FPhi} shifted by $q$ bounds $e_R$ and $e_R'$ by
$A_je^{-a_qe^{2|x|}}$ outside $[-R,R]$; integrating gives $\|e^{|x|}e_R\|_2^2\le\frac{A_0^2}{2a_q}e^{-2a_qe^{2R}}$ and
$\|e_R'\|_2^2\le\frac{A_1^2+c_\chi^2A_0^2}{a_q}e^{-2a_qe^{2R}}$. Finally
$\D(u)\le\frac23\|u'\|_2^2+\frac{32}3\|u\|_2^2$ for $u\in H^1$, from $\|\Delta_tu\|_2\le\min(t\|u'\|_2,2\|u\|_2)$ and
$a(t)\le\frac4{3t}$ on $(0,1]$, $a(t)\le\frac43e^{-t/2}$ on $[1,\infty)$.
\end{proof}

\begin{lemma}[independence of translates]\label{lem:indep}
Let $f\in L^1(\R)$ be continuous with $\int f\ne0$. If $p_1,\dots,p_k$ are distinct and
$\sum_\ell d_\ell f(p_\ell-q)=0$ for all $q\in\mathbb Q$, then all $d_\ell=0$.
\end{lemma}
\begin{proof}
By continuity the identity holds for all real $q$; Fourier transform in $q$ gives $\widehat f(-\xi)\sum_\ell d_\ell e^{-i\xi p_\ell}=0$.
Since $\widehat f(0)=\int f\ne0$, the exponential polynomial vanishes on an interval around $0$; its derivatives of order
$0,\dots,k-1$ at $0$ give $\sum_\ell d_\ell p_\ell^j=0$, and the Vandermonde determinant $\prod_{i<j}(p_j-p_i)\ne0$ forces $d=0$.
\end{proof}

The theorem is best stated in the generality in which it is proved, and in which it is formalised (\Cref{app:lean}): nothing
about the Weil form enters except the two facts that $f_0$ is a positive integrable function and that its cutoffs have vanishing
form value.

\begin{theorem}[no positive finite-stencil minorant]\label{thm:nominorant}
Let $f>0$ be continuous and integrable on $\R$, let $\chi_R$ be measurable cutoffs with $\chi_R=1$ on $[-R,R]$, let
$\Q_0:(\R\to\C)\to\R$ be any functional with
\[
 \Q_0\bigl(f\,\chi_R\,f(\cdot-q)/f\bigr)\longrightarrow0\quad(R\to\infty)\qquad\text{for every }q\in\mathbb Q .
\]
Let $(Ss)(x)=\sum_{\ell=1}^kc_\ell\,s(x+\tau_\ell)$ with distinct shifts $\tau_\ell$ and a nonzero real coefficient
vector, and let $W\ge0$ be measurable. If $\int W\,|S(\chi_Rf(\cdot-q)/f)|^2\,dx\le\Q_0(f\chi_Rf(\cdot-q)/f)$ for all $q\in\mathbb Q$ and
$R\ge1$, then $W=0$ almost everywhere. No bound on $k$, on the diameter of the stencil, or the condition $\sum c_\ell=0$ is needed.
\end{theorem}
\begin{proof}
Fix $q\in\mathbb Q$ and write $s_R=\chi_Rf(\cdot-q)/f$, $r=f(\cdot-q)/f$. The hypotheses give $\int W|Ss_R|^2\to0$ along $R=1,2,\dots$.
For fixed $x$ all $k$ sample points $x+\tau_\ell$ lie in $[-R,R]$ once $R\ge|x|+\max|\tau_\ell|$, where $Ss_R(x)=Sr(x)$; the integrands
$W|Ss_R|^2$ are nonnegative and converge pointwise to $W|Sr|^2$, so Fatou's lemma gives $\int W|Sr|^2=0$. Hence for each rational $q$
there is a null set $N_q$ off which $W(x)=0$ or $Sr(x)=0$. Off the null set $\bigcup_qN_q$, every $x$ with $W(x)>0$ satisfies
$\sum_\ell\frac{c_\ell}{f(x+\tau_\ell)}f(x+\tau_\ell-q)=0$ for all $q\in\mathbb Q$; \Cref{lem:indep} with $p_\ell=x+\tau_\ell$ and
$d_\ell=c_\ell/f(x+\tau_\ell)$ (and $\int f>0$) forces $c=0$, a contradiction. Hence $W=0$ a.e.
\end{proof}

\begin{corollary}[Weil form]\label{cor:nominorantWeil}
Let $S$, $W$ be as in \Cref{thm:nominorant}. If $\int W|Ss|^2\,dx\le\Q(f_0s)$ for all $s\in C_c^\infty(\R;\C)$, then $W=0$ a.e.
\end{corollary}
\begin{proof}
Take $f=f_0$, $\Q_0=\Q$, and smooth $\chi_R$ as in \Cref{lem:budget}; the hypothesis of the theorem on the functional is \eqref{eq:budget},
and the minorant hypothesis is a special case of the one assumed here, since $s_{q,R}\in C_c^\infty$.
\end{proof}

The noncompact ratio $r_q$ is never inserted into the hypothesis; only the compact $s_{q,R}$ are, and the limit is taken
on the certificate side. Nothing about the sign of $\Q$ or the location of zeros enters.

\begin{corollary}[no sum-of-squares representation]\label{cor:noCSS}
There is no countable family of finite stencils $S_j$ and weights $W_j\ge0$ with
$\Q(f_0s)=\sum_j\int W_j|S_js|^2dx$ for all $s\in C_c^\infty$.
\end{corollary}
\begin{proof}
Each nonnegative summand is a minorant, hence zero by \Cref{cor:nominorantWeil}, so $\Q(f_0\,\cdot)$ would vanish identically.
But it does not: for a nonnegative bump $g$ with $\|g\|_2=1$ supported in an interval of length
$\ell\le\min\{\tfrac12,\tfrac12\log2,\exp(-e^{1/2}(c_A+1))\}$ the prime correlations vanish, the pole term is positive,
$\|\Delta_tg\|_2^2=2$ for $\ell\le t\le1$ and $a(t)\ge e^{-1/2}/(2t)$ there, so $\Q(g)\ge e^{-1/2}\log(1/\ell)-c_A\ge1$.
\end{proof}

\begin{remark}
The same conclusion holds for families $\int_J\!\int W(j,x)|S_js(x)|^2dx\,d\mu(j)$ over a $\sigma$-finite parameter space with jointly
measurable data and finitely many sample points per $j$: by Tonelli the total is a nonnegative integral over $(j,x)$, Fatou and the countable
set of rational $q$ give $W(j,x)=0$ or $S_jr_q(x)=0$ for a.e.\ $(j,x)$ simultaneously for all $q$, and \Cref{lem:indep} applies pointwise.
We do not use this extension and leave the measure-theoretic details to the reader.
\end{remark}

Quantitatively, if a proposed summand has $W>0$ on a set of positive measure, there are a rational $q$, a bounded set $E$
and $\eta=\int_EW|Sr_q|^2>0$ such that for $R$ large (explicitly $e^{2R}>(2a_q)^{-1}\log(1+2C_{\Xf}C_q^2/\eta)$ and $R$
covering $E+\{\tau_\ell\}$)
\begin{equation}\label{eq:certkill}
 \Q(f_0s_{q,R})-\int W|Ss_{q,R}|^2\le-\eta/2<0 .
\end{equation}
This is a strict negative \emph{certificate margin}; it is not a negative value of $\Q$.

\begin{example}[the obstruction is not a sign statement]\label{ex:control}
Let $F(x)=e^{-(x-1)^2}+e^{-(x+1)^2}$, $\omega=\pi/2$, and $\Q_*(g)=|\int g(x)e^{-i\omega x}dx|^2\ge0$. Since
$\widehat F(\omega)=2\cos\omega\cdot\widehat{e^{-x^2}}(\omega)=0$, all translates of $F$ lie in the radical of $\Q_*$, $F>0$ is
integrable, and the argument above applies verbatim: the manifestly nonnegative, nonzero form $\Q_*(F\,\cdot)$ admits no
positive finite-stencil minorant either. Its positive factor is a global integral, not a finite sample.
\end{example}

\begin{remark}[finite graphs]
For a real symmetric matrix $M$ with $M\mathbf 1=0$, positive semidefiniteness is equivalent to $M=\sum\lambda_iu_iu_i^{\!\top}$
with $\lambda_i\ge0$ and $u_i\perp\mathbf 1$ (spectral theorem). The continuum statement fails because the ratio form has far
more null functions than constants, and by \Cref{lem:indep} their restrictions to any finite set of points span all sample
values. Two-point edge squares are an even narrower cone: $(1,-2,1)(1,-2,1)^{\!\top}$ is PSD with zero row sums but has a
positive $(1,3)$ entry, so it is not a positive combination of edge squares.
\end{remark}

\subsection{Two further remarks on proof shapes}\label{sec:trade}

\begin{proposition}[no uniform margin on a dense family]\label{prop:trade}
Let $\mathcal F\subset\Xf$ be a family whose linear span is dense in $\Xf$. Then there is no $\delta>0$ with
$\Q(g)\ge\delta\|g\|_{\Xf}^2$ for all $g$ in the span of $\mathcal F$.
\end{proposition}
\begin{proof}
$\Q$ and $\|\cdot\|_{\Xf}^2$ are continuous on $\Xf$ (\Cref{sec:setup}), so the inequality would extend to all of $\Xf$; but
$\Q(f_0)=0$ and $\|f_0\|_{\Xf}>0$ (\Cref{prop:radical}).
\end{proof}
The consequence for proof design is that positivity of $\Q$ cannot be established by proving a fixed $\Xf$-margin on a dense family
of tests and passing to the limit: any margin that survives densification is a coercivity statement, and the form is not coercive. The
finite windows show the same thing quantitatively: $\lambda_1(m,N)$ decays exponentially in $m$ (\Cref{sec:windows}).

\begin{remark}[pointwise majorisation]\label{rem:schur}
A second proof shape is to bound the negative part of the kernel of $\Q(f_0\,\cdot)$ pointwise by a positive majorant built from the positive
part. We record, without a theorem, why we do not expect this to work: the positive density $b_+(t)\sim\frac1{2t}$ is not integrable at $t=0$
against the autocorrelation, while its usable energy against a fixed test is only $\int_0^\delta b_+E_s\,dt=\frac{\delta^2}4\int f_0^2|s'|^2+o(\delta^2)$;
the positive part therefore cannot be transported to the negative region $t>t_0$ by any $s$-independent pointwise comparison. A precise
impossibility statement would have to specify the class of majorants, and we do not pursue it here.
\end{remark}
